\ExerciseSection

\begin{enumerate}
\def\labelenumi{\arabic{enumi})}
\tightlist
\item
  Let \(\varphi\) be the canonical homomorphism of \(\mathbf{R}^n\) onto \(\mathbf{T}^n\) and let \(u\) be a linear mapping of \(\mathbf{R}^m\) into \(\mathbf{R}^n\) . For \(\varphi \circ u\) to be a strict morphism of \(\mathbf{R}^m\) into \(\mathbf{T}^n\) it is necessary and sufficient that either:
\end{enumerate}

\begin{enumerate}
\def\labelenumi{\alph{enumi})}
\item
  \(\overline{u}^1 (\mathbf{Z}^n)\) is of rank \(m\) ; or
\item
  \(u(\mathbf{R}^m)\cap \mathbf{Z}^n\) is of rank equal to the dimension of \(u(\mathbf{R}^m)\) .
\end{enumerate}

\begin{enumerate}
\def\labelenumi{\arabic{enumi})}
\setcounter{enumi}{1}
\item
  Show that the only continuous homomorphism of the additive topological group \(Q_{p}\) of p-adic numbers (Chapter III, § 6, Exercises 23 et seq.) into the additive group R is the zero mapping.
\item
  Every \(p\) -adic number \(x\) (Exercise 2) can be written in the form \(\sum_{-\infty}^{+\infty} \alpha_k p^k\) where the \(\alpha_k\) are rational integers, those of index \(k < 0\) being zero for all but a finite number of values of the index \(k\) . Moreover, if \(\sum_{-\infty}^{+\infty} \alpha_k p^k = \sum_{-\infty}^{+\infty} \beta_k p^k\) , the rational numbers \(\sum_{-\infty}^{-1} \alpha_k p^k\) and \(\sum_{-\infty}^{-1} \beta_k p^k\) are congruent mod 1. Let \(\varphi_p(x)\) be the class mod 1 of the rational number \(\sum_{-\infty}^{-1} \alpha_k p^k\) .\\
\end{enumerate}

\begin{enumerate}
\def\labelenumi{\alph{enumi})}
\tightlist
\item
  Show that \(\varphi_p\) is a continuous homomorphism of the topological group \(\mathbf{Q}_p\) into \(\mathbf{T}\) , and that for every continuous homomorphism \(f\) of \(\mathbf{Q}_p\) into \(\mathbf{T}\) there exists a unique \(a \in \mathbf{Q}_p\) such that \(f(x) = \varphi_p(ax)\) for all \(x \in \mathbf{Q}_p\) . {[}Note that the knowledge of the elements \(f(p^k)\) for \(k \leqslant 0\) determines \(f\) uniquely; show that each of these elements is the class mod 1 of a fraction whose denominator is a power of \(p\) ; hence show that there exists \(a \in \mathbf{Q}_p\) such that \(f(p^k) = \varphi_p(ap^k)\) for all \(k \leqslant 0\) .{]}\\
\item
  Show that, if \(\varphi: \mathbf{R} \to \mathbf{T}\) is the canonical homomorphism, then \(\varphi(x) = \sum_{p} \varphi_p(x)\) for all \(x \in \mathbf{Q}\) , the sum being over all prime numbers \(p\) (observe that all but a finite number of terms of this sum will be zero).
\end{enumerate}

¶ 4) a) Let G be a locally compact abelian group and let H be a subgroup of G isomorphic to \(\mathbf{R}^{p} \times \mathbf{T}^{q}\) where \(p, q\) are integers \(\geqslant 0\) ; suppose also that G/H is isomorphic to either R or T. Show that G is isomorphic to the product of H and a closed subgroup L isomorphic to G/H. {[}Apply Chapter V, § 3, Exercise 7 d), to obtain a continuous homomorphism \(f\) of R into G such that \(f(\mathbf{R}) \subset \mathbf{H}\) ; hence construct another continuous homomorphism \(g: \mathbf{R} \to \mathbf{G}\) such that \(f(t) - g(t) \in \mathbf{H}\) for all \(t \in \mathbf{R}\) and \(g(\mathbf{R}) \cap \mathbf{H} = \{e\}\) . Note that \(f(\mathbf{H})\) is a proper closed subgroup of R and that, for every element \(z \neq e\) in H, there exists a continuous homomorphism \(u\) of R into H such that \(u(1) = z\) .{]}

\begin{enumerate}
\def\labelenumi{\alph{enumi})}
\setcounter{enumi}{1}
\tightlist
\item
  Let G be a Hausdorff abelian topological group, let H be a closed subgroup of G isomorphic to \(R^{p} \times T^{q}\) , and suppose that there exists a surjective continuous homomorphism \(\varphi: R \to G/H\) . Show that there exists a continuous homomorphism \(f: R \to G\) such that \(f(\mathbf{R}) \cap \mathbf{H} = \{e\}\) and \(\mathbf{G} = \mathbf{H}. f(\mathbf{R})\) . {[}Reduce to case a) by considering the least upper bound on G of the given topology and the topology for which a funda-
\end{enumerate}

mental system of neighbourhoods of o is formed by the \(\overline{p}(\varphi(I))\) , where \(p:G\to G/H\) is the canonical mapping and I runs through a fundamental system of neighbourhoods of o in R; show that G is locally compact in this new topology (cf.~Chapter III, § 4, Exercise 10 e).{]}

\begin{enumerate}
\def\labelenumi{\alph{enumi})}
\setcounter{enumi}{2}
\tightlist
\item
  Give an example where the restriction of \(p\) to \(f(\mathbf{R})\) is not bicontinuous {[}cf.~§ 1, Exercise 2f){]}.
\end{enumerate}

\begin{enumerate}
\def\labelenumi{\arabic{enumi})}
\setcounter{enumi}{4}
\item
  Let G be a connected topological group and let H be a compact normal subgroup of G with no arbitrarily small subgroups (Chapter III, § 2, Exercise 30). Show that H is contained in the centre of G. (Observe that for s close to e in G, the image of H under \(x \rightarrow sxs^{-1}\) is arbitrarily close to e.)
\item
  Let \(\mathbf{G}\) be a topological group and let \(\mathbf{H}\) be a closed normal subgroup of \(\mathbf{G}\) , isomorphic to \(\mathbf{R}^n (n\geqslant 1)\) and such that \(\mathbf{G} / \mathbf{H}\) is abelian.
\end{enumerate}

\begin{enumerate}
\def\labelenumi{\alph{enumi})}
\item
  Suppose, moreover, that H contains no connected subgroup normal in G except for H itself and \(\{e\}\) , and that H is not contained in the centre of G. Given \(x_{0} \in G\) which does not commute with every element of H, show that, for every \(v \in H\) , there exists a unique \(u \in H\) such that \(x_{0}^{-1}ux_{0}u^{-1} = v\) . (Observe that in the additive notation in H the equation becomes \(u - x_{0}^{-1}ux_{0} = -v\) and that the continuous endomorphism \(u \to u - x_{0}^{-1}ux_{0}\) of H is a linear mapping of \(R^{n}\) into itself. Show that its kernel is a normal subgroup of G.) Furthermore, if \(f(u)\) is the unique \(u \in H\) such that \(x_{0}^{-1}ux_{0}u^{-1} = v\) , then f is a bicontinuous automorphism of H.
\item
  Under the hypotheses of a), show that the continuous mapping \(g: y \to (f(x_{0}^{-1}yx_{0}y^{-1}))^{-1}y\) of G into itself has as image the subgroup L of G consisting of all elements which commute with \(x_{0}\) and that the relation \(g(y) = g(y')\) is equivalent to \(yy'^{-1} \in H\) . Deduce that there exists a continuous bijection \(h: G/H \to L\) whose inverse is the restriction to L of the canonical mapping \(p: G \to G/H\) , and deduce that G is isomorphic to the topological semi-direct product of H and L.
\end{enumerate}

\begin{enumerate}
\def\labelenumi{\arabic{enumi})}
\setcounter{enumi}{6}
\tightlist
\item
  Let G be a topological group, H a normal subgroup isomorphic to \(R^{p} \times T^{q}\) ( \(p \geqslant 0\) , \(q \geqslant 0\) ), and suppose that there exists a surjective continuous homomorphism \(\varphi: R \to G/H\) . Show that there exists a continuous homomorphism \(f: R \to G\) such that \(f(R) \cap H = \{e\}\) and \(G = H\) . \(f(R)\) . {[}Using Exercise 4 b) and Chapter V, § 3, Exercise 8, reduce to the case where G is not commutative and then, using Exercise 5, to the case q = 0; now use induction on p, together with Exercise 6.{]}
\end{enumerate}

